%=====================================================================
\chapter{Systematic Literature Review and Systematic Mapping}\label{ch:slr}
%=====================================================================
\locator{2}{Part~2 --- Finding and Synthesising Knowledge}

\begin{outcomes}
By the end of this chapter you will be able to:
\begin{tight}
  \item \textbf{Distinguish} a systematic review, a systematic mapping study
    and a narrative review, and choose between them.
  \item \textbf{Write} a review protocol before conducting the search.
  \item \textbf{Construct} a search string from PICOC and adapt it to the
    syntax of four digital libraries.
  \item \textbf{Screen} studies against explicit criteria and report inter-rater
    agreement.
  \item \textbf{Appraise} study quality with an instrument appropriate to the
    designs under review.
  \item \textbf{Synthesise} findings without resorting to vote counting.
  \item \textbf{Report} the whole process as a PRISMA 2020 flow.
\end{tight}
\end{outcomes}

\begin{prereq}
\chref{questions} for PICOC, which drives the search string, and for the
question types --- a review's questions have types too. \chref{contribution}
for why a review that does not synthesise is not a contribution.
\end{prereq}

\begin{keyterms}
systematic review $\bullet$ systematic mapping study $\bullet$ protocol
$\bullet$ primary study $\bullet$ search string $\bullet$ digital library
$\bullet$ inclusion and exclusion criteria $\bullet$ screening $\bullet$
inter-rater agreement $\bullet$ Cohen's kappa $\bullet$ quality assessment
$\bullet$ data extraction $\bullet$ synthesis $\bullet$ vote counting
$\bullet$ publication bias $\bullet$ PRISMA
\end{keyterms}

\section{A Review Is a Study}

The narrative review --- read widely, form an impression, write it up --- is
how most people first encounter the literature, and it has a fatal property:
nobody can check it. Its conclusions depend on which papers the author happened
to find, which is unrecorded, and on which they found convincing, which is
unexamined. Two competent people reviewing the same field this way routinely
reach different conclusions, and there is no way to adjudicate.

A \term{systematic review} treats the literature as data. It has a protocol
written in advance, a search that another person could re-run, explicit
criteria applied by more than one screener, an appraisal stage, and a synthesis
that is answerable to what was found rather than to what was remembered. It is
a study whose subjects happen to be papers, and every methodological concern in
this book applies to it: bias, reliability, validity, reproducibility.

\begin{center}
\small
\begin{tabular}{@{}L{2.6cm}L{3.7cm}L{3.7cm}L{3.4cm}@{}}
\toprule
 & \textbf{Systematic review} & \textbf{Mapping study} & \textbf{Narrative review} \\
\midrule
Question    & Specific: what is the effect of X on Y?
            & Broad: what has been researched, and where are the gaps?
            & Whatever the author chooses \\
Search      & Exhaustive within a stated scope
            & Exhaustive, usually broader
            & Unstated \\
Appraisal   & Quality assessed, may exclude
            & Usually classification only
            & None \\
Synthesis   & Answers the question
            & Categorises the field; counts by facet
            & Argues a position \\
Output      & An answer with confidence in it
            & A map, and identified gaps
            & An orientation \\
Use it when & The question is answerable from existing evidence
            & The area is unfamiliar or immature
            & Writing a thesis background section, honestly labelled \\
\bottomrule
\end{tabular}
\end{center}

\begin{conceptbox}[Which one does your synopsis need?]
Almost always a mapping study first, then a systematic review if the map shows
enough primary studies to answer a specific question.

\medskip
A scholar entering an unfamiliar area cannot yet write the specific question a
systematic review requires. The map produces it: once you can see that 60
papers address microservice testing but only four evaluate anything in
production, the specific question writes itself --- and you have the evidence
that it is a genuine gap, which is exactly what \chref{contribution} demands of
a delta.
\end{conceptbox}

\begin{paperbox}[The two guideline documents to work from]
Barbara Kitchenham and Stuart Charters, \emph{Guidelines for Performing
Systematic Literature Reviews in Software Engineering}, EBSE-2007-01,
2007~\cite{kitchenham2007guidelines}.\oa{}

\medskip
Matthew Page and colleagues, \emph{The PRISMA 2020 statement}, BMJ 372:n71,
2021~\cite{page2021prisma}.\oa{}

\medskip
\textbf{Why they matter.} Kitchenham and Charters is the procedure the software
engineering community actually follows; PRISMA 2020 is the reporting standard
reviewers and editors increasingly require, and its flow diagram is expected in
any serious review. They are complementary: the first tells you what to do, the
second what to report.

\medskip
\textbf{What to extract.} From Kitchenham, the three-phase structure --- plan,
conduct, report --- and the protocol contents. From PRISMA, the 27-item
checklist and the flow diagram, and from PRISMA-S~\cite{rethlefsen2021prisma}
the requirement to report the full search string for \emph{every} database, not
just one.
\end{paperbox}

\section{The Protocol Comes First}

The protocol is written, reviewed by your supervisor, and fixed \emph{before}
the search runs. This is not bureaucracy: it is the same device as
pre-registration (\chref{prereg}), and it exists to stop you deciding what
counts as relevant after you have seen which papers support the conclusion you
were hoping for.

\begin{templatebox}[C3 --- Systematic review protocol]
\textbf{1. Rationale.} Why is this review needed now? What existing reviews
are there, and why are they insufficient?
\filllines{2}

\noindent\textbf{2. Review questions.} Numbered. Each must be answerable from
published primary studies.
\filllines{3}

\noindent\textbf{3. PICOC.} From template C2.
\filllines{2}

\noindent\textbf{4. Search strategy.} Databases, date range, and the exact
string per database (attach as an appendix).
\filllines{3}

\noindent\textbf{5. Inclusion criteria.} What must a study have to be included?
\filllines{2}

\noindent\textbf{6. Exclusion criteria.} Stated separately, not as negations of
the above.
\filllines{2}

\noindent\textbf{7. Screening procedure.} Who screens, at which stages, how
disagreement is resolved, and how agreement is measured.
\filllines{2}

\noindent\textbf{8. Quality assessment instrument.} Which, and whether low
quality excludes or is analysed as a moderator.
\filllines{2}

\noindent\textbf{9. Data to extract.} The fields of template C4.
\filllines{1}

\noindent\textbf{10. Synthesis method.} Narrative, thematic, or meta-analysis
--- decided now, not after seeing the data.
\filllines{2}
\end{templatebox}

\section{Building the Search}

\subsection{From PICOC to a string}

Take the PICOC elements from template C2. For each element that will constrain
the search --- usually P, I and C --- list synonyms, abbreviations, spelling
variants and older terminology. Join synonyms with \texttt{OR}, join elements
with \texttt{AND}.

\begin{worked}[A search string, built and then repaired]
\textbf{Question.} What techniques are used to test microservice-based systems,
and which have been evaluated in industrial settings?

\medskip
\textbf{Attempt 1.}
\begin{lstlisting}[language=,basicstyle=\ttfamily\footnotesize]
microservice AND testing
\end{lstlisting}
Returns far too much, and misses more. ``Microservices'' (plural) is not
matched by every engine; ``micro-service'' and ``micro service'' are used in
older papers; ``testing'' misses ``test'', ``validation'', ``verification''.

\medskip
\textbf{Attempt 2.}
\begin{lstlisting}[language=,basicstyle=\ttfamily\footnotesize]
("microservice*" OR "micro-service*" OR "micro service*")
AND
("test*" OR "verification" OR "validation" OR "quality assurance")
\end{lstlisting}
Better, but ``validation'' pulls in a large machine-learning literature about
model validation that has nothing to do with the question.

\medskip
\textbf{Attempt 3 --- add the context element.}
\begin{lstlisting}[language=,basicstyle=\ttfamily\footnotesize]
("microservice*" OR "micro-service*" OR "micro service*")
AND
("test*" OR "verification" OR "integration testing"
 OR "end-to-end testing" OR "contract testing")
AND
("software" OR "system*" OR "architecture")
\end{lstlisting}

\medskip
\textbf{Then calibrate.} Assemble a \emph{quasi-gold standard}: five to ten
papers you already know must be included. Run the string. If it does not
return all of them, it is not yet good enough --- inspect each miss and find
the term that would have caught it. Record this calibration in the protocol; it
is the only evidence a reviewer has that your search was adequate.

\medskip
All counts and terms here are illustrative.
\end{worked}

\begin{pitfall}[Every database has different syntax, and the differences bite]
\begin{tight}
  \item \textbf{IEEE Xplore} supports nested parentheses and \texttt{*}
    wildcards, but caps command-search string length --- long strings must be
    split and the results unioned.
  \item \textbf{ACM Digital Library} changed its search interface and syntax;
    strings written for the old one silently return different results.
  \item \textbf{Scopus} and \textbf{Web of Science} use field codes
    (\texttt{TITLE-ABS-KEY}, \texttt{TS=}). Omitting the field code searches
    full text and returns thousands of irrelevant hits.
  \item \textbf{SpringerLink} handles Boolean operators inconsistently across
    content types.
  \item \textbf{Google Scholar} should not be used as a primary source: results
    are not reproducible between users or dates, and there is no export of a
    complete result set. Use it for snowballing (\chref{greylit}), not for the
    systematic search.
\end{tight}
\medskip
Record the exact string, the database, the fields searched, the date, and the
number of hits --- separately for each. PRISMA-S requires
it~\cite{rethlefsen2021prisma}, and without it your review is not reproducible.
\end{pitfall}

\section{Screening}

\begin{center}
\small
\begin{tabular}{@{}L{2.6cm}L{4.6cm}L{7.0cm}@{}}
\toprule
\textbf{Stage} & \textbf{What is read} & \textbf{Typical retention} \\
\midrule
Deduplication & Metadata & Overlap between IEEE, ACM and Scopus is
                           substantial; a reference manager does this \\
Title screen  & Title only & Fast, high recall, deliberately generous \\
Abstract screen & Title and abstract & Where most exclusion happens \\
Full-text screen & The whole paper & Slowest; the only stage where you can
                                     apply criteria about method or data \\
Snowballing   & Reference lists and citing papers & Adds papers the string
                                                    missed (\chref{greylit}) \\
\bottomrule
\end{tabular}
\end{center}

\subsection{Criteria that can actually be applied}

An inclusion criterion is usable when two people applying it to the same paper
reach the same answer. ``Studies of high quality'' is not usable. ``Studies
reporting at least one quantitative outcome measure with a sample size stated''
is.

\begin{conceptbox}[Measuring agreement, and what the number means]
Have a second screener --- your supervisor, a fellow scholar --- independently
screen a sample, typically 10--20\% of records. Report Cohen's $\kappa$, which
corrects for agreement expected by chance:
\[
\kappa \;=\; \frac{p_o - p_e}{1 - p_e}
\]
where $p_o$ is observed agreement and $p_e$ the agreement expected if both
screeners were guessing at their observed marginal rates.

\medskip
Landis and Koch's conventional bands~\cite{landis1977measurement} --- 0.41--0.60
moderate, 0.61--0.80 substantial, above 0.81 almost perfect --- are widely
quoted and were explicitly described by their authors as arbitrary. Treat them
as a rough signal, not a threshold to be passed.

\medskip
A low $\kappa$ is useful information: it means the criteria are ambiguous. The
fix is to repair the criteria and re-screen, not to discuss individual papers
until you agree. \chref{qualitative} covers $\kappa$ and its alternatives in
more depth.
\end{conceptbox}

\section{Quality Appraisal}

Not every published study deserves equal weight. Appraisal asks whether a
study's design supports its conclusions --- not whether you like its findings.

The instrument must match the designs under review. A checklist built for
randomised trials cannot appraise a case study. In computing, a practical
approach is a short list of design-appropriate questions, each scored 1 /
0.5 / 0:

\begin{center}
\small
\begin{tabular}{@{}L{0.6cm}L{13.8cm}@{}}
\toprule
Q1 & Are the aims and research questions clearly stated? \\
Q2 & Is the design appropriate to those aims (\chref{questions})? \\
Q3 & Is the sample or subject selection described and justified? \\
Q4 & Are the measures defined well enough to be reproduced
     (\chref{measurement})? \\
Q5 & Is the analysis appropriate, and are assumptions checked
     (\chref{inference})? \\
Q6 & Are threats to validity identified and discussed (\chref{validity})? \\
Q7 & Are the data or artefacts available (\chref{reproducibility})? \\
\bottomrule
\end{tabular}
\end{center}

\begin{alertbox}[Decide the role of quality scores in the protocol]
There are two defensible uses, and you must choose in advance:
\begin{tight}
  \item \textbf{As an exclusion threshold} --- studies below a stated score are
    excluded. Simple, but discards evidence and the threshold is arbitrary.
  \item \textbf{As a moderator} --- all studies are included, and you check
    whether high- and low-quality studies give different answers. More
    informative, and it turns a methodological worry into a finding.
\end{tight}
\medskip
What is not defensible is scoring quality after synthesis and using it to
explain away studies that disagreed with your conclusion.
\end{alertbox}

\section{Extraction and Synthesis}

\begin{templatebox}[C4 --- Data extraction form]
One row per included study. Pilot it on five papers before extracting all of
them; you will always find fields you cannot populate.

\medskip
\begin{tabular}{@{}L{4.0cm}L{9.5cm}@{}}
Study ID / citation      & \fillline{9.2cm} \\[6pt]
Year, venue, venue type  & \fillline{9.2cm} \\[6pt]
Research question(s)     & \fillline{9.2cm} \\[6pt]
Study type               & \fillline{9.2cm} \\[6pt]
Population / subjects, N & \fillline{9.2cm} \\[6pt]
Intervention             & \fillline{9.2cm} \\[6pt]
Comparison               & \fillline{9.2cm} \\[6pt]
Outcome measures         & \fillline{9.2cm} \\[6pt]
Effect size / result     & \fillline{9.2cm} \\[6pt]
Context (industrial / academic / simulated) & \fillline{9.2cm} \\[6pt]
Quality score            & \fillline{9.2cm} \\[6pt]
Threats acknowledged     & \fillline{9.2cm} \\[6pt]
Artefacts available      & \fillline{9.2cm} \\[6pt]
Relevance to RQ1 / RQ2 / RQ3 & \fillline{9.2cm} \\[6pt]
\end{tabular}
\end{templatebox}

\subsection{Three ways to synthesise, and one to avoid}

\textbf{Narrative synthesis} groups studies and describes patterns in prose. It
is the default in software engineering, because primary studies are rarely
commensurable enough for anything stronger. Done well it explains \emph{why}
studies disagree --- different populations, different operationalisations,
different baselines. Done badly it is a list.

\textbf{Thematic synthesis} codes findings across studies and builds themes,
using the machinery of \chref{qualitative}. Appropriate when the primary
studies are themselves qualitative.

\textbf{Meta-analysis} pools effect sizes statistically, weighting by precision.
It requires that studies measure comparable constructs with reported effect
sizes and variances, which in computing is uncommon --- but where it is
possible, it is far stronger than anything else. The \texttt{metafor} package
in R is the standard open-source tool.

\begin{pitfall}[Vote counting is not synthesis]
``Twelve studies found a positive effect, five found none, therefore the effect
is real.'' This reasoning is wrong, and it is everywhere.

\medskip
It ignores sample size, so a study of 12 participants counts as much as one of
1,200. It ignores effect size, so a trivial effect counts as much as a large
one. And it is biased against exactly the truth it claims to establish: several
underpowered studies finding nothing are perfectly compatible with a real
effect, and a set of small studies finding an effect is compatible with
publication bias (\chref{prereg}) and no effect at all.

\medskip
If effect sizes are reported, pool them. If they are not, say so, describe the
pattern, and treat the absence of pooled evidence as itself a finding about the
state of the field.
\end{pitfall}

\section{Reporting: the PRISMA Flow}

\armfig{
\begin{tikzpicture}[
  bx/.style={draw=primary!60, fill=primary!6, rounded corners=2pt,
             text width=4.6cm, align=center, font=\footnotesize,
             inner sep=4pt, minimum height=0.95cm},
  ex/.style={draw=accent!60, fill=accent!6, rounded corners=2pt,
             text width=5.0cm, align=left, font=\footnotesize,
             inner sep=4pt, minimum height=0.95cm},
  ph/.style={font=\scriptsize\bfseries, text=secondary, rotate=90},
  ar/.style={-{Stealth[length=2.2mm]}, draw=secondary, semithick}]

  \node[bx] (id)  at (0,0)    {Records identified\\ from databases ($n=\underline{\hspace{9mm}}$)};
  \node[bx] (dd)  at (0,-1.65){Records after duplicates removed ($n=\underline{\hspace{9mm}}$)};
  \node[bx] (sc)  at (0,-3.30){Records screened on title and abstract ($n=\underline{\hspace{9mm}}$)};
  \node[bx] (ft)  at (0,-4.95){Full texts assessed for eligibility ($n=\underline{\hspace{9mm}}$)};
  \node[bx] (sn)  at (0,-6.60){Additional records from snowballing ($n=\underline{\hspace{9mm}}$)};
  \node[bx, fill=secondary!12, draw=secondary!70]
            (in)  at (0,-8.25){Studies included in synthesis ($n=\underline{\hspace{9mm}}$)};

  \node[ex] (e1) at (6.6,-1.65) {Duplicates removed ($n=\underline{\hspace{7mm}}$)};
  \node[ex] (e2) at (6.6,-3.30) {Records excluded, with reason ($n=\underline{\hspace{7mm}}$)};
  \node[ex] (e3) at (6.6,-4.95) {Full texts excluded, \emph{by reason} ($n=\underline{\hspace{7mm}}$)};

  \foreach \a/\b in {id/dd, dd/sc, sc/ft, ft/sn, sn/in} {\draw[ar] (\a) -- (\b);}
  \draw[ar] (dd) -- (e1);
  \draw[ar] (sc) -- (e2);
  \draw[ar] (ft) -- (e3);

  \node[ph] at (-3.5,-0.8)  {IDENTIFICATION};
  \node[ph] at (-3.5,-4.1)  {SCREENING};
  \node[ph] at (-3.5,-8.25) {INCLUDED};
\end{tikzpicture}}
{The PRISMA 2020 flow, as a form to fill in. The single most scrutinised box is
the third exclusion node: reporting \emph{how many} full texts were excluded is
routine, reporting them \emph{broken down by reason} is what distinguishes a
review a reader can trust. Reproduced in fill-in form as template
C5.}{prismaflow}

\begin{conceptbox}[What the numbers reveal to a reader]
An experienced reviewer reads your flow diagram diagnostically.

\medskip
A very high title-screen exclusion rate --- 4,000 down to 60 --- suggests the
search string was too broad and probably imprecise. A very low one suggests it
was too narrow and may have missed relevant work. A large snowballing yield
relative to the database yield says the search string was inadequate, and
honest reviews say so. And full-text exclusions with no stated reasons are the
clearest possible signal that criteria were applied after the fact.
\end{conceptbox}

\begin{regbox}[Why this chapter matters for your synopsis --- PhD \S13, MS \S8]
The synopsis is defended before the Departmental Research Committee, and for
PhD scholars additionally evaluated by an external referee from a national
university, chosen from a panel of six experts supplied by your supervisor
(\reg{PhD Regs 2024}{\S13(ii)}).

\medskip
The reviewer's first question is always whether the gap is real. A protocol, a
search string and a PRISMA flow answer that question with evidence. An
assertion that ``little work exists in this area'' invites the referee to name
three papers that do, and that is a bad afternoon.
\end{regbox}

\begin{traditions}{reviewing the literature}
\tradEMP{Pool effect sizes where the primary studies permit it; where they do
not, report why not, and treat the incommensurability as a finding about the
field's measurement practices.}
\tradDSR{Review artefacts and design knowledge rather than effects: what has
been built, for which problem class, evaluated how convincingly.}
\tradQUAL{Thematic or meta-ethnographic synthesis, building second-order
constructs across studies. Aims at interpretation, not at a pooled estimate.}
\tradFORM{Review results and their assumptions: which bounds are known, under
which models, and where the open problems sit. Exhaustiveness matters more than
sampling.}
\end{traditions}

\begin{lab}[Run the first two stages for real]
Using template C3, write a protocol for a review on your own topic. Then:

\begin{tightnum}
  \item Assemble a quasi-gold standard of five papers you know must be included.
  \item Build a search string from your PICOC. Run it on IEEE Xplore, ACM DL
    and Scopus. Record the exact string, fields, date and hit count for each.
  \item Check recall against your gold standard. Repair and re-run. Record how
    many iterations it took --- reviewers of your synopsis will ask.
  \item Deduplicate in Zotero or JabRef and screen the first 100 titles.
  \item Have a colleague screen the same 100. Compute $\kappa$.
\end{tightnum}

\medskip
If $\kappa$ is below about 0.6, rewrite the criteria rather than arguing about
the papers. This is portfolio piece P2.
\end{lab}

\begin{checkpoint}
\begin{tightnum}
  \item Why must the protocol be fixed before the search is run? Name the
    specific bias this prevents.
  \item Your string returns 12,000 records. Give three distinct repairs and say
    what each risks.
  \item A review reports ``143 full texts assessed, 38 included''. What is
    missing, and why does its absence matter?
  \item Explain why nine studies finding an effect and four finding none is not
    evidence that the effect is real.
\end{tightnum}
\end{checkpoint}

\begin{chaptersummary}
\begin{tight}
  \item A systematic review is a study whose subjects are papers; it has a
    protocol, a reproducible search and an appraisal stage, and it can be
    checked.
  \item Map first when the area is unfamiliar; the map generates the specific
    question a review needs.
  \item Write the protocol before searching. It is pre-registration for a
    review.
  \item Build the string from PICOC, calibrate it against a quasi-gold
    standard, and record string, fields, date and hits per database.
  \item Screen with criteria two people can apply identically; report $\kappa$
    and repair the criteria, not the disagreement.
  \item Decide the role of quality scores in advance --- exclusion threshold or
    moderator.
  \item Synthesise; do not vote count. If effect sizes exist, pool them; if not,
    explain the disagreement.
  \item Report as a PRISMA 2020 flow, with full-text exclusions broken down by
    reason.
\end{tight}
\end{chaptersummary}

\begin{reviewq}
  \item Systematic reviews were imported into software engineering from
    evidence-based medicine. Identify two features of computing research that
    make the transfer imperfect, and say how the guidelines accommodate --- or
    fail to accommodate --- each.
  \item A reviewer excludes all studies scoring below 5 on a 7-point quality
    instrument, and the remaining studies agree. A critic argues this
    manufactured the agreement. Assess the criticism and describe an analysis
    that would settle it.
  \item Google Scholar has the best coverage of any source and is excluded from
    systematic searches. Explain the tension, and propose a defensible policy
    for its use.
  \item Publication bias means the literature you can find is not the research
    that was done. What can a reviewer in computing actually do about this,
    given that trial registries do not exist for our field?
\end{reviewq}
